problem:  
Let \( E_{p,q} = SU(3)//S^1_{p,q} \) be an Eschenburg space given by the free circle action  
\[
z \cdot A = \operatorname{diag}(z^{p_1},z^{p_2},z^{p_3})\,A\,\operatorname{diag}(\overline{z}^{\,q_1},\overline{z}^{\,q_2},\overline{z}^{\,q_3}),
\]
where \(p,q\in\mathbb{Z}^3\) and \(\sum_i p_i=\sum_i q_i\).  
Let \(L\) be a **left-invariant** metric on \(SU(3)\) that is **right-invariant** under \(S^1_{p,q}\); denote by \(g_{p,q}(L)\) the induced **submersion metric** on \(E_{p,q}\).  

Suppose the identity component \(\mathrm{Isom}_0(E_{p,q}, g_{p,q}(L))\) of the isometry group acts with **cohomogeneity one** on \(E_{p,q}\).  

**Research problem:**  
Is it possible for \(g_{p,q}(L)\) to have **nonnegative sectional curvature everywhere** while containing *no* region of strictly positive curvature under this cohomogeneity‑one hypothesis?  
Equivalently, does the existence of a cohomogeneity‑one, right‑invariant submersion metric on an Eschenburg space **force positivity of curvature somewhere**?  

In a broader sense, does cohomogeneity one serve as a geometric rigidity condition that excludes entirely flat directions for \(g_{p,q}(L)\) among submersion metrics on Eschenburg spaces?

Why is it a "good" problem:  
This reformulated problem is mathematically precise and structurally sound. It focuses on the **interaction between isometric cohomogeneity-one actions and curvature positivity** within the tightly controlled family of **submersion metrics** on Eschenburg spaces—metrics that descend from left-invariant, right-\(S^1_{p,q}\)-invariant metrics on \(SU(3)\).  

The question probes whether intermediate symmetry (cohomogeneity one) enforces geometric rigidity strong enough to guarantee the appearance of positive curvature somewhere, even when only nonnegative curvature is assumed a priori. This operates at the intersection of  

- **transformation group theory**, examining how the dimension and orbit structure of the isometry group constrain geometry;  
- **Riemannian submersion and curvature formulae**, which relate horizontal and vertical distributions in the biquotient framework; and  
- **curvature rigidity phenomena**, central to the modern study of positively curved manifolds.  

A resolution would clarify whether symmetry itself can induce local curvature positivity in exotic 7‑manifolds, potentially illuminating the transition between homogeneous, nearly homogeneous, and nonhomogeneous curvature behavior. The statement is clean—"Does cohomogeneity one force some positivity on Eschenburg spaces?"—yet any serious attempt to answer it requires integrating deep tools from differential geometry, topology of biquotients, and group actions, embodying the "narrow entrance, profound depth" aesthetic of an excellent research problem.